\section{Experimental protocols and additional results}
\label{app:experiment}

All task-success counts in this appendix refer to simulation. Primary comparisons pair
(task, initial state), without pinned policy sampling; a fixed episode changes
failure to success, and a broken episode changes success to failure. Healthy runs measure nominal task competence. They
are distinguished from faulted-frozen controls and from healthy runs with adaptation
enabled. Early diagnostic tables without individual paired outcomes are reported using
marginal counts only.

\subsection{Calibration, identification, and evaluation protocols}
\label{app:protocols}

\begin{center}\small
\setlength{\tabcolsep}{4pt}
\begin{tabular}{>{\raggedright\arraybackslash}p{.17\linewidth}>{\raggedright\arraybackslash}p{.25\linewidth}>{\raggedright\arraybackslash}p{.22\linewidth}>{\raggedright\arraybackslash}p{.25\linewidth}}
\toprule
Interface & Healthy calibration & Fault identification & Evaluation correction\\
\midrule
LIBERO Panda & Three spatial episodes, 285 steps, with $\pi_{0.5}$; current plus six delayed commands & Within the evaluation episode & Episode-local updates; mask selected for the fault family\\[3pt]
ALOHA & Eight episodes; joint-position predictor & One preceding faulted episode for the held protocol & Estimate held during 40 evaluation episodes\\[3pt]
GR1 & Six episodes; arm-position predictor & Reach windows, steps 50--200, in three preceding faulted episodes & Median estimate held during 30 evaluation episodes per condition\\
\bottomrule
\end{tabular}
\end{center}
LIBERO reuses the predictor and sensitivity across suites, OpenVLA-OFT, and GR00T N1.7.
ALOHA and GR1 each receive robot-specific calibration. Their held protocols include the
listed identification episodes in addition to evaluation. Continuous-adaptation runs
are separate conditions. Thus policy transfer on one interface and adaptation on a
new interface use different calibration budgets.

\paragraph{Implementation provenance.}
A registered provenance pass recovered, for all nine cohorts, the code
path, gate convention, clipping order, $K$, law constants, and calibration
identity, archived as an implementation map with per-run configuration
records. Every published cohort subtracts the correction outside the
network; the native output-bias edit exists only in oracle and diagnostic
scripts. The gate is a leakage deadzone on the LIBERO and ALOHA cohorts and
a full hold on GR1; it never fired in the four headline samples or ALOHA's
identification episodes, and did fire in the joint-level cells, healthy
controls, and the deadzone-$.03$ ablation row. $K=6$ in every cohort. Known
residual caveats are retained: the shipped sensitivity replay overran its
source episode by five commands (the held-out calibration does not share
this), two of twenty spatial episodes share a (task, init) pair with the
plant log, constants were tuned on the same initial-state band, and the
early unpaired four-suite tables used per-suite calibration where the paired
headline uses the spatial one. The held-out cells answer the overlap
caveats; no historical row is retrospectively reassigned to a cleaner
update law.

The LIBERO primary condition estimates six motion coordinates and applies correction
to the three rotation coordinates. Translation and rotation masks in the fault-family
sweep were selected using faulted separation diagnostics. Healthy data fit the predictor
and local sensitivity; fault-family development observations informed mask selection.
A preregistered healthy-only gate (three-sigma phantom bands from twenty
healthy estimate-only episodes) reproduces the rotation mask without faulted
data; no automatic unknown-support selection beyond that gate is claimed. A model fit on healthy behavior also requires checking its error when the
policy visits new contact and state regimes.
The healthy predictor budget is tabulated above. LIBERO sensitivity uses central
probes of $\pm.02$; total probe/reset counts, the fault-development selection budget,
and complete task/seed identities for the final evaluations are not documented.
Thus the entire development procedure cannot be described as healthy-data-only.

\paragraph{Architecture illustration and online information flow.}
\label{app:figure-record}
Figure~\ref{fig:framework} depicts only the evaluated prediction-only
information path. At step $k$, the frozen VLA supplies $a_k$ and the available
snapshot $\widehat f_{\tau_k}$ is masked by $D$ before realization and clipping.
The resulting $u_k$ is recorded at the modeled interface; the physical fault
then produces $u_k^{\rm eff}=u_k+f_k$. The post-command response $y_k$ and the
issued-command history enter the calibrated FIR predictor. Their residual is
transformed by $S^{-1}$ and healthy centering, then the declared attenuation
or innovation rule updates or holds the estimate. The resulting
$\widehat f_{k+1}$ can only affect a later command through the snapshot rule,
so the diagram contains no same-step algebraic loop.

The mask $D$ selects corrected coordinates and remains distinct from the
residual normalizer $P_r$. Predictor coefficients $H_\ell$, sensitivity $S$,
healthy center $b_{\rm cal}$, limits, and $D$ are fixed before deployment.
The figure is a mechanism diagram rather than empirical evidence: it contains
no task-success cohort, nominal replay, task reward, fault truth, VLA update,
geodesic signal, or online metric computation.

\subsection{Healthy centering and physical error measurements}
\label{app:centering}
The model uses $z_k=S^{-1}e_k^{\rm res}-b_{\rm cal}$, with $b_{\rm cal}$ in the
same coordinates as the estimated fault. Historical descriptions also call
the baseline a settled healthy estimate. These definitions need not agree.
For a constant healthy inverse-sensitivity observation $b_0$, constant
attenuation $s_0$, and inactive projection, the uncentered attenuated filter
settles at $s_0b_0$. Subtracting that quantity before the same attenuation
leaves a target $s_0(1-s_0)b_0$, which is generally nonzero. With varying
gains, $\mathbb E[s_kz_k]\ne\mathbb E[z_k]$ in general. Recomputing the gain
after centering changes this illustrative expression again. Exact averaging
quantity, units, subtraction location, and whether the normalizer is centered
must therefore be recovered before assigning healthy drift to a physical fault.
The existing counts are retained under their documented protocols.

For a simulator diagnostic, injected truth can reveal $w_k=z_k-f_k$,
$\|D(\fhat_k-f_k)\|$, the uncorrected component $\|(I-D)f_k\|$, and the
applied snapshot age. Truth is for evaluation, not deployment feedback.
Decoder realization error must be isolated using base and edited actions at
an identical observation and decoder draw; comparing different rollouts
confounds realization with state change. Log projection and command clipping
separately. Phase-resolved residual error, persistence and incremental-motion
prediction baselines, and time until a prespecified execution tolerance is
sustained can connect recovery to the theoretical budget. Those raw records
are absent from the supplied experiment summaries, so no new VLA onset plot,
physical-time recovery curve, or empirical metric estimate is reconstructed.
The constructed example in Appendix~\ref{app:contraction-example} supplies
such curves for its explicitly specified servo only.
For execution-error measurements, replay the deployment sequence
$a_k$ through the nominal physical robot--controller model
with matched initial conditions and exogenous inputs. The comparison
$\xi_{k}^{\rm nom}$ uses those recorded commands; a healthy VLA rerun usually produces
a different sequence and estimates task competence instead. A model-based
replay also needs a valid model-error allowance. Controller state, physical
delays, and the execution metric must describe the same subsystem.

\subsection{Statistical interpretation and repeated evaluations}
\label{app:statistics}

For $b$ fixed and $c$ broken episodes, the two-sided exact McNemar test is
\begin{equation}
 p_{\mathrm{McN}}=\min\left\{1,
 2^{1-(b+c)}\sum_{j=0}^{\min(b,c)}\binom{b+c}{j}\right\},
 \qquad p_{\mathrm{McN}}=1\text{ if }b+c=0.
\end{equation}
The primary four-suite comparison has $b=51,c=1$, giving
$p=2.354\times10^{-14}$ and an episode-weighted gain of $50/120=41.7$ percentage
points. The unweighted mean of the four suite improvements is $44.4$ points.
Pointwise 95\% Wilson intervals are $[16.7,31.7]\%$ for $28/120$ and
$[56.1,72.9]\%$ for $78/120$; these describe the two marginal rates, not their
paired difference. The four primary suite comparisons remain significant under either
Holm or Bonferroni correction within that family of four tests.

For the paired effect, define $d_i=Y_i^{\mathrm{corrected}}-Y_i^{\mathrm{frozen}}
\in\{-1,0,1\}$ and $\widehat\Delta=(b-c)/n$. Under independent, identically
distributed episode pairs, an estimated standard error is
\begin{equation}
 \widehat{\mathrm{se}}(\widehat\Delta)
 =\sqrt{\frac{b+c-n\widehat\Delta^2}{n(n-1)}}.
 \label{eq:paired-se}
\end{equation}
This follows from $\sum_i d_i^2=b+c$ and the unbiased sample variance of $d_i$.
For $n=120,b=51,c=1$, the standard error is $.04672$; the normal approximation
$\widehat\Delta\pm1.96\,\widehat{\mathrm{se}}$ gives $[.3251,.5082]$.
This is an approximate episode-iid 95\% interval, not an exact interval or a
task-clustered confidence statement. The discordant counts also imply 27 pairs
successful in both arms and 41 unsuccessful in both. Individual episode identities
are unnecessary for this count-based calculation but are necessary to estimate
task clustering or construct a task-level bootstrap.

A task-clustered bootstrap ($B=20{,}000$, resampling the ten tasks per
suite) gives pooled 95\% intervals of $[+36,+65]$ episodes under the shipped
calibration and $[+20,+49]$ under the held-out one, both excluding zero.
Per suite, cluster-level sign-flip $p$ under the shipped calibration is
$.031$, $.016$, $.061$, $.007$ (spatial, object, goal, long-horizon) and
under the held-out one $.009$, $.081$, $.191$, $.127$: three of four suites
and one of four, respectively, against four of four at episode level. The
pooled effect is robust to task clustering; individual suites at $n=20$--$40$
are not.
Other condition-level tests are exploratory and nominal.
Repeated evaluation of one recorded $\pi_{0.5}$ condition gives $13/20$,
$13/20$, and $15/20$, and the friction rerun's identical twenty scenarios
scored $8$, $4$, and $6$ corrected across three runs, illustrating the
variation in these small samples. Nonsignificance
against healthy performance does not establish equivalence or noninferiority.
Matched healthy controls on the exact four primary samples now exist:
frozen 116/120 and 118/120 across two reruns (spatial 19--20, object 20,
goal 39--40, long-horizon 38), and with the law running 117/120 and 118/120
($b{=}4,c{=}3$ and $b{=}2,c{=}2$; $p=1.0$), so the corrected rates are 67\%
(shipped) and 58\% (held-out) of health on the same scenarios. The separate
\code{libero\_90} competence check gives $6/20$ healthy successes and is
excluded from the recovery headline.

Two severity conditions were repeated at $n=40$ using the reported initial-state set
45--48. These are separate evaluations:
\begin{center}\small
\begin{tabular}{lrrrrr}
\toprule
Fault & Frozen & Corrected & Fixed & Broken & Exact $p$\\
\midrule
Rotation $0.05$ & $27/40$ & $39/40$ & 12 & 0 & $4.883\times10^{-4}$\\
Translation $0.10$ & $30/40$ & $38/40$ & 9 & 1 & $0.0215$\\
\bottomrule
\end{tabular}
\end{center}
The rotation-$0.05$ condition is significant in this repetition, whereas the original
$n=20$ cell has $p=.0625$. The translation repetition includes one regression.
Zero regressions in the original severity cells therefore describe the observed
sample rather than a general guarantee of harmless adaptation.

\subsection{Calibration and the measured residual}
\label{app:calibration}

The LIBERO predictor uses the current command and six delayed commands. Its reported
calibration-fit scores are:
\begin{center}\small
\begin{tabular}{lrrrrrr}
\toprule
Predictor $R^2$ & $dx$ & $dy$ & $dz$ & $dr_x$ & $dr_y$ & $dr_z$\\
\midrule
Static & .885 & .961 & .978 & .416 & .091 & .077\\
FIR, chart subtraction & .975 & .982 & .982 & .489 & .109 & .168\\
FIR, relative rotation & .975 & .982 & .982 & .522 & .336 & .972\\
\bottomrule
\end{tabular}
\end{center}
Relative orientation uses $\operatorname{Log}(R_{t+1}R_t^{\mathsf T})^\vee$ in a
consistent frame. A numerical check over 2,000 random increments of magnitude
$0.05$ rad reports maximum error $1.4\times10^{-14}$, versus $6.28$ for subtraction
of axis--angle coordinates. The latter can cross a chart discontinuity. These fit
and representation checks motivate the residual construction; they are not error
bounds on unseen contact modes. For absolute-position outputs on ALOHA and
GR1, a persistence predictor and incremental-motion error are necessary
comparisons: high position $R^2$ alone does not identify command-induced
motion, physical faults, or a contraction rate.

Per-axis central probes of $\pm0.02$ in the experiment's action coordinates give
\begin{equation}
S=\begin{bmatrix}
 .297& .021&-.112&-.029& .081&-.006\\
 .008& .272& .023&-.034& .016& .003\\
 .029& .032& .126&-.037& .012&-.003\\
-.003&-.001&-.001& .253&-.004& .002\\
 .013& .001&-.001&-.007& .276& .000\\
 .001& .004&-.003& .011& .005& .244
\end{bmatrix}.
\end{equation}
The reported condition number is $3.0$, and about 26\% of the sum of absolute entries
is off the diagonal. Repeating probes at $\pm0.05$ changes diagonal entries by ratios
$(.99,.92,1.14,.99,.91,.99)$ and gives reported condition number $2.2$. This supports an approximately stable single-axis sensitivity over the
probed \emph{magnitude} range. Configuration is different: re-probing at
held-out initial state 25 lowers the translation diagonal by 19--21\% on
$y,z$ (rotation within 8\%), and at a third initial state the $z$ entry more
than doubles ($.126\to.262$) while the condition number falls from $3.0$ to
$1.29$ -- two of three registered stability bands failed. A single-state
probe can carry a factor-of-two error in one axis, and the shipped matrix's
small $z$ entry is a property of its probe state, not of the plant.

Simultaneous offsets produce larger deviations from superposition. The ratio of the
measured uniform-fault response to the sum of single-axis responses is $0.77$--$1.04$
at magnitudes $0.01$ and $0.02$. At magnitude $0.05$, translation ratios depend on
probe direction: $x$ gives $0.07$ and $0.78$, and $z$ gives $0.50$ and $0.27$.
Contact and command limits are possible contributors; these measurements do not
separate their individual effects. Corresponding separation measurements are:
\begin{center}\small
\begin{tabular}{lrr}
\toprule
Injected fault / measured channel & Estimated separation & Fraction of true offset\\
\midrule
Uniform six-axis $0.05$ / $x$ & $-0.0130$ & Wrong sign\\
Single-axis $0.05$ / $x$ & $+0.0244$ & $49\%$\\
Single-axis $0.02$ / $x$ & $+0.0158$ & $79\%$\\
Single-axis $0.02$ / $r_x$ & $+0.0201$ & $101\%$\\
\bottomrule
\end{tabular}
\end{center}
Translation is therefore identifiable in some probed regimes even though the
uniform-fault evaluation uses rotation-only correction. Local sensitivity, model
error, and the realized fault family jointly determine useful correction directions.

\subsection{Severity map and structured offsets}
\label{app:severity}
\begin{figure}[H]
\centering\includegraphics[width=\linewidth]{fig_map.pdf}
\caption{Reported severity sweep, LIBERO spatial: $n=20$ paired episodes per condition. Bars show success; whiskers show marginal 95\% Wilson intervals, not paired-effect intervals. Translation faults use translation correction; rotation and uniform faults use rotation correction. Supports are family-informed, not autonomously selected.}
\label{fig:map}
\end{figure}

All cells in Figure~\ref{fig:map} use 20 paired spatial-suite episodes. Translation
faults use translation correction; rotation and uniform faults use rotation correction.
\begin{center}\small
\begin{tabular}{llrrrrr}
\toprule
Family & Magnitude & Frozen & Corrected & Fixed & Broken & Exact $p$\\
\midrule
Translation & $.05$ & $18/20$ & $18/20$ & 0 & 0 & $1$\\
 & $.10$ & $13/20$ & $19/20$ & 6 & 0 & $.03125$\\
 & $.15$ & $4/20$ & $19/20$ & 15 & 0 & $6.104\times10^{-5}$\\
Rotation & $.05$ & $13/20$ & $18/20$ & 5 & 0 & $.0625$\\
 & $.10$ & $0/20$ & $17/20$ & 17 & 0 & $1.526\times10^{-5}$\\
 & $.15$ & $0/20$ & $9/20$ & 9 & 0 & $.003906$\\
Uniform six-axis & $.05$ & $8/20$ & $18/20$ & 10 & 0 & $.001953$\\
 & $.10$ & $2/20$ & $11/20$ & 9 & 0 & $.003906$\\
 & $.15$ & $0/20$ & $2/20$ & 2 & 0 & $.5$\\
\bottomrule
\end{tabular}
\end{center}
The uniform-$0.05$ cell repeats the primary spatial condition. At uniform offset
$0.15$, the remaining translation disturbance and operation outside the calibrated
regime are consistent with the limited recovery. This comparison does not uniquely
identify a failure mechanism.
Figure~\ref{fig:map} displays separate marginal Wilson intervals for each arm.
For $k$ successes among $n=20$, $\hat p=k/n$, and $z=1.959964$, the limits are
\begin{equation}
 \frac{\hat p+z^2/(2n)\;\pm\;
 z\sqrt{\hat p(1-\hat p)/n+z^2/(4n^2)}}{1+z^2/n}.
\end{equation}
These whiskers describe uncertainty in individual success rates under an
episode-binomial model. They do not encode within-pair dependence or replace
Eq.~\eqref{eq:paired-se} for the primary paired difference.

For mixed-sign rotation offset $(+.06,-.06,+.03)$, success improves $2/12\to10/12$.
The separation estimate is $(+.053,-.038,+.029)$: all signs agree with the fault,
with magnitude errors $11.7\%$, $36.7\%$, and $3.3\%$. Zero-fault translation
channels read $(-.003,-.024,-.006)$, showing nonuniform off-channel leakage.
Discordant counts are unavailable for this exploratory condition, so no paired
significance test is assigned.

\subsection{Ablations of the implemented offset estimator}
\label{app:ablation}

Each row changes one setting relative to the LIBERO spatial reference: uniform
$+0.05$ fault, rotation-only correction, $\gamma=0.08$, $\rho=0.15$, residual deadzone
$0.008$, and a predictor fitted on three healthy episodes (285 steps). The reference
repeats the primary spatial condition; its row is the archived headline file
itself, not a fresh execution of the ablation runner. Endpoint estimates are last-50-step averages;
healthy-subtracted separation estimates use a different measurement protocol.
\begin{center}\small
\begin{tabular}{lrrrrrl}
\toprule
Setting & Frozen & Corr. & Fixed & Broken & Exact $p$ & Endpoint $(r_x,r_y,r_z)$\\
\midrule
Reference & $8/20$ & $18/20$ & 10 & 0 & .00195 & (.044,.020,.044)\\
$\gamma=.02$ & $8/20$ & $17/20$ & 9 & 0 & .00391 & (.036,.016,.036)\\
$\gamma=.32$ & $7/20$ & $18/20$ & 11 & 0 & .00098 & (.044,.019,.044)\\
$\rho=.05$ & $11/20$ & $14/20$ & 5 & 2 & .453 & (.027,.011,.027)\\
$\rho=.50$ & $8/20$ & $20/20$ & 12 & 0 & .00049 & (.048,.022,.049)\\
Deadzone $0$ & $8/20$ & $19/20$ & 11 & 0 & .00098 & (.044,.018,.044)\\
Deadzone $.03$ & $8/20$ & $17/20$ & 9 & 0 & .00391 & (.031,.016,.032)\\
One healthy episode & $10/20$ & $17/20$ & 8 & 1 & .0391 & (.027,.008,.044)\\
Two healthy episodes & $8/20$ & $17/20$ & 10 & 1 & .0117 & (.045,.021,.043)\\
\bottomrule
\end{tabular}
\end{center}
Success stays at $17/20$ or higher over the tested $16\times$ range of $\gamma$.
Frozen totals differ across settings, so these within-row paired comparisons do not
form a directly paired between-method ablation. Common task/initial-state/seed
cohorts are needed to attribute differences to the changed setting alone.
A smaller normalizer reduces estimated offsets and introduces two regressions in its
sample. Under constant residual and inactive projection, the retention factor is
$s=(1+\|r\|^2/\rho^2)^{-1}$. At the representative residual norm $0.034$, it
is $0.951$ for $\rho=.15$ and $0.684$ for $\rho=.05$. The latter differs from the
approximately 55\% observed retention on two axes. Residual variation, calibration
bias, and finite-time behavior must also be accounted for to predict that magnitude.
At $\gamma=.08$, 36 active updates reduce the transient toward a constant attenuated
target by at least 95\%. This does not mean a 95\% reduction in error to the true
fault. Attenuation could limit harmful correction when the predictor is inaccurate,
but no matched attenuation-versus-innovation experiment establishes that mechanism.

\subsection{Time-varying offsets}
\label{app:time}

The following evaluations use amplitude $0.10$:
\begin{center}\small
\begin{tabular}{lrrl}
\toprule
Profile & Frozen & Corrected & Reported exact $p$\\
\midrule
Translation ramp over 60 steps & $12/20$ & $18/20$ & .031\\
Rotation oscillation, nonzero mean & $13/20$ & $20/20$ & .016\\
Intermittent rotation & $12/20$ & $19/20$ & .039\\
\bottomrule
\end{tabular}
\end{center}
The tests for the ramp and intermittent profiles are retained as reported; their
marginal counts do not determine the discordant pairs. Success on every corrected
oscillatory episode determines seven fixed and zero broken pairs, giving
$p=.015625$. Waveform RMS errors are 45--79\% of fault amplitude. Useful correction
therefore occurs despite imperfect waveform tracking in these conditions. A zero-mean
oscillatory condition leaves frozen performance essentially intact and provides no
recovery evidence.

\subsection{Gain estimation and the integral baseline}
\label{app:gain}

Replacing the instantaneous-command regressor by the FIR-weighted history reduces
the healthy gain-estimation phantom from $0.181$ to $0.020$, and increases
$x$-channel separation from $1.4$ to $18.1$ healthy-phantom standard deviations.
The corresponding healthy success comparisons are $10/10\to9/10$ and
$10/10\to10/10$. This diagnostic directly supports time alignment of the gain
regressor.

Translation gain $g=.30$ gives $1/20\to17/20$ with reported $p=3.1\times10^{-5}$.
The estimates $(-.723,-.699,-.717)$ have absolute errors $(.023,.001,.017)$ relative
to $\beta^\star=-.700$, or $(3.3,.1,2.4)\%$ of $|\beta^\star|$.
At $g=.20$, success improves $0/20\to17/20$, although some estimates saturate at the
projection boundary. At $g=.50$, the comparison is $15/20\to18/20$ with reported
$p=.25$. Rotation channels lacking the required excitation in the evaluated
trajectories do not support gain-identification claims.

The raw-error integral baseline is $\fhat\leftarrow\fhat+k_i(y_t-a_t)$, evaluated
under translation offset $0.15$:
\begin{center}\small
\begin{tabular}{lrr}
\toprule
Integral gain & Frozen & Integral correction\\
\midrule
.001 & $4/20$ & $5/20$\\
.005 & $3/20$ & $7/20$\\
.020 & $5/20$ & $0/20$\\
.050 & $2/20$ & $1/20$\\
.150 & $4/20$ & $1/20$\\
\midrule
Model-based method & $4/20$ & $19/20$\\
\bottomrule
\end{tabular}
\end{center}
The best observed baseline success is 35\%, compared with 95\% for
model-based correction, but frozen totals vary across this gain sweep. A
seed-pinned paired rerun on one scenario set settles it: method $18/20$
against integral $3/20$, 15 fixed and 0 broken, $p=6.1\times10^{-5}$. The
estimates explain the gap -- method $112/56/125\%$ of the fault on $x/y/z$
against integral $18/41/163\%$, the integral's $x$ never arriving and its $z$
overshooting. The comparison supports calibrating the command response; the
separate static-observer cell isolates FIR memory, and a calibrated
disturbance observer remains uncompared.

\subsection{Native-bias reachability diagnostic}
\label{app:oracle}

An early $\pi_{0.5}$ experiment applies a computed six-vector to
\code{action\_out\_proj/bias} under uniform $+0.05$ offset, scaled by $k$.
Each condition uses 15 episodes:
\begin{center}\small
\begin{tabular}{lr}
\toprule
Intervention & Successes\\
\midrule
Faulted, $k=-1$ & $0/15$\\
Faulted, $k=0$ (no edit) & $7/15$\\
Faulted, $k=.5$ & $13/15$\\
Faulted, $k=1$ (computed oracle) & $15/15$\\
Faulted, $k=1.5$ & $3/15$\\
Healthy, $k=1$ & $1/15$\\
\bottomrule
\end{tabular}
\end{center}
This diagnostic establishes a useful correction at the native bias site in
that setup and substantial sensitivity to its sign and scale. It uses a
different intervention and sample from the primary rotation-only comparison.
A registered measurement later characterized the edit at fixed observations:
deterministic and linear there (relative linear mismatch $.035$, quadratic
part at most 8\%), but with decoder gain varying 27--41\% across
observations, so a single mean-Jacobian bias vector realizes an intended
$\pm.05$ correction with median relative error 43\% (95th percentile 1.9,
maximum 4.5 times the intended magnitude). This table therefore neither
establishes exact external-action cancellation nor supplies an oracle upper
bound for every evaluation protocol; the same-interface known-fault oracle
in the main text does. Individual paired outcomes are unavailable.

\subsection{Faults below the operational-space controller}
\label{app:joint}

These interventions act on the simulated Panda elbow below the controller.
Correction uses translation channels with the healthy phantom subtracted --
the one cohort that subtracts a settled attenuated estimate, measured
open-loop and applied closed-loop. The historical friction and lock cells
ran with model mutation on a cached environment and without a guaranteed
common initial scene; a registered rerun repeats them with enforced model
restoration and scenario reset.
\begin{center}\small
\begin{tabular}{lrrrrr}
\toprule
Intervention & Frozen & Corrected & Fixed & Broken & Exact $p$\\
\midrule
Torque bias $5$ N\,m (historical) & $20/40$ & $32/40$ & 15 & 3 & .00754\\
Friction $+4$ (historical) & $20/20$ & $20/20$ & 0 & 0 & 1\\
Friction $+10$ (historical) & $11/20$ & $14/20$ & 6 & 3 & .508\\
Friction $+20$ (rerun, $n{=}40$) & $0/40$ & $13/40$ & 13 & 0 & $2.44\cdot10^{-4}$\\
Lock within $\pm.05$ rad (rerun) & $0/20$ & $0/20$ & 0 & 0 & 1\\
Healthy elbow, law running (rerun) & $19/20$ & $20/20$ & 1 & 0 & 1\\
\bottomrule
\end{tabular}
\end{center}
The frozen damage reproduces under the corrected protocol (friction $1/20$
and $0/40$ against the historical $0/20$; lock $0/20$), so the harm was not
a protocol artefact. The friction repair did not reproduce at $n{=}20$
($4/20$ against a registered keep-threshold of five); a preregistered
$n{=}40$ extension resolved it at $0/40\to13/40$ with no regressions,
replacing the historical 40\% with a measured 33\% from a floor of zero.
The lock cell reproduces exactly, estimate included (final $x,z$ estimates
$-.232,-.189$ against the recorded $-.23,-.16$): identified, corrected, not
repaired. On a healthy elbow the law holds a translation phantom of up to
$.167$ and costs nothing. The torque comparison, on the historical protocol,
also includes three regressions. The lock yields a residual separation of
15--23 healthy-phantom standard deviations but no recovery: strong detection
does not ensure sufficient correction authority through the tested action
interface; other controllers or replanned motions remain outside this
comparison.

\subsection{Additional backbones on the LIBERO interface}
\label{app:transferdetails}

All rows use the unchanged LIBERO calibration. Offset rows use the indicated fault
family; gain $.20$ acts on translation.
\begin{center}\small
\begin{tabular}{lllrrrr}
\toprule
Policy & Suite & Condition & Frozen & Corrected & Fixed & Broken\\
\midrule
OpenVLA-OFT & Spatial & Healthy & $20/20$ & $19/20$ & 0 & 1\\
 & Spatial & Rotation $.10$ & $0/20$ & $11/20$ & 11 & 0\\
 & Spatial & Translation $.15$ & $0/20$ & $14/20$ & 14 & 0\\
 & Spatial & Gain $.20$ & $0/20$ & $17/20$ & 17 & 0\\
 & Object & Rotation $.10$ & $0/20$ & $10/20$ & 10 & 0\\
 & Object & Translation $.15$ & $0/20$ & $15/20$ & 15 & 0\\
 & Object & Gain $.20$ & $0/20$ & $19/20$ & 19 & 0\\
\midrule
GR00T N1.7 & Spatial & Healthy & $18/20$ & $18/20$ & 1 & 1\\
 & Spatial & Rotation $.10$ & $0/20$ & $14/20$ & 14 & 0\\
 & Spatial & Translation $.15$ & $2/20$ & $15/20$ & 13 & 0\\
 & Spatial & Gain $.20$ & $0/20$ & $19/20$ & 19 & 0\\
\bottomrule
\end{tabular}
\end{center}
For OFT object-suite rotation, translation, and gain conditions, exact paired tests
give $p=.00195$, $6.104\times10^{-5}$, and $3.815\times10^{-6}$, respectively.
The healthy rows quantify regressions that aggregate rates can obscure: N1.7 retains
the same total while changing two individual outcomes. Successful calibration reuse
shows transfer over the evaluated trajectories; model error can still depend on the
policy through state, contact, and excitation distributions.

\subsection{ALOHA: identification and held correction}

ALOHA transfer-cube uses 14 absolute joint targets at 50 Hz and 300-step episodes.
Calibration on eight healthy episodes gives joint-position fit scores
$R^2=.989$--$1.000$ on the twelve arm joints (the uncorrected right gripper
fits at $.77$); the arm-block sensitivity condition number is $1.17$, and
the previously quoted $7.3$ belongs to the full matrix including that
gripper's $.145$ diagonal. At left-arm offset
$+.02$ rad:
\begin{center}\small
\begin{tabular}{lrrr}
\toprule
Condition & Correction variation & Mean residual & Successes\\
\midrule
Faulted frozen & --- & $2.12$ cm & $0/20$\\
Continuous adaptation (attenuation law) & $.43$ cm & $.17$ cm & $0/20$\\
Static $.019$ rad correction & $0$ & $.11$ cm & $5/20$\\
Prior identification, then hold & $0$ & $.07$ cm & $15/40$\\
Faulted frozen, matched $n=40$ & --- & --- & $0/40$\\
Healthy, matched $n=40$ & --- & --- & $17/40$\\
Healthy, adaptation active & --- & --- & $14/40$\\
\bottomrule
\end{tabular}
\end{center}
The held comparison has 15 fixed and zero broken episodes, giving
$p=6.104\times10^{-5}$. Its one preceding identification episode produces estimates
$.019$--$.020$ rad, or 95--101\% of the injected offset. The healthy active-law
comparison has reported $p=.58$; discordant counts are unavailable.

Separation at offsets $.05$ and $.02$ rad is 85--91\% and 94--106\% of the fault,
respectively, with leakage at most $.007$ rad on other joints. The $.05$ rad offset
displaces the gripper by $5.3$ cm. These observations show useful identification, while the
continuous-adaptation condition illustrates that small mean residual alone
does not ensure task success. A registered matched-initialization comparison
then separated the confound: from the same stored estimate, the innovation
law updating throughout repairs $14/40$ against the held $15/40$ ($p=1.0$)
while moving $.30$ cm within episodes, and the attenuation law repairs
$0/40$ while moving $.40$ cm -- its estimate is pulled to $.016$ rad by step
twenty (the attenuated fixed point plus deadzone leakage once the fault is
cancelled), a $.42$ cm standing error during the grasp. The failure under
continued adaptation is the attenuation law's bias, not updating as such. LIBERO's $.019$ figure is the mean per-step change of a velocity-offset
correction (about $.019$ cm per commanded unit step, a realized arm-motion
change of $.0044$ cm), while ALOHA's $.43$ cm is the median within-episode
range (SD $.13$ cm) of a position offset: the two are not on one scale, and
the margin comparison across interfaces is qualitative.
High absolute-position $R^2$ can also arise from slowly varying joint
positions, and the registered diagnostic confirms it: the FIR's position
$R^2$ is within $.001$ of a persistence predictor's. On position
\emph{increments} the held-out FIR reaches $.53$ (ALOHA left arm) and $.38$
(GR1 plate task) against persistence near zero, and $-1.77$ on the GR1 tray
task, where it is worse than persistence; the correction relies on the
plant's steady-state gain, within a few percent of one on every corrected
joint except the tray-task wrist ($.91$, $.94$).

\subsection{GR1: fault magnitude and task dependence}

The Fourier GR1 upper-body simulator uses GR00T N1.5, distinct from the N1.7 policy on
LIBERO. There are 29 absolute joint targets at 20 Hz. Six healthy episodes fit the plate-task arm predictor with $R^2\ge.998$;
the tray-task cells use a separate ten-episode plant fitting at
$.994$--$.9997$. Measured arm sensitivity is approximately $.99I$. The residual norm excludes the twelve unmodeled hand joints. Identification
uses reach-window measurements, steps 50--200, across three preceding faulted
episodes; the median estimate is held during evaluation.
\begin{center}\small
\begin{tabular}{llrrrrrl}
\toprule
Task & Offset & Healthy & Frozen & Corr. & Fixed & Broken & Exact $p$\\
\midrule
Plate & $.05$ & $19/30$ & $7/30$ & $24/30$ & 17 & 0 & $1.526\times10^{-5}$\\
Plate & $.10$ & $22/30$ & $1/30$ & $19/30$ & 18 & 0 & $7.629\times10^{-6}$\\
Plate & $.20$ & $25/30$ & $0/30$ & $21/30$ & 21 & 0 & $9.537\times10^{-7}$\\
Tray & $.10$ & $17/30$ & $1/30$ & $13/30$ & 12 & 0 & $4.883\times10^{-4}$\\
\bottomrule
\end{tabular}
\end{center}
Offsets are in radians and apply to all seven working-arm joints. The four rows
contain 120 evaluation pairs with zero observed regressions; the plate-$.10$ and
tray-$.10$ rows repeat Table~\ref{tab:transfer}. Held estimates span 90--108\%,
93--104\%, and 94--99\% of the true plate offsets, and 73--126\% on the tray task.
Continuous adaptation with the attenuation law on plate-$.10$ gives
$8/30$, versus $19/30$ when held; in the registered matched-initialization
rerun on one episode set, the innovation law continued gives $17/30$ against
$22/30$ held, and the attenuation law $1/30$ -- on this robot updating
itself costs five episodes, and holding is what the small task margin
requires. Direct discordant pairs for the historical continuous-versus-held
and healthy-versus-corrected comparisons are unavailable, so their
significance is not reconstructed.

\subsection{Identification traces and timing scope}
\label{app:identificationtraces}
\begin{figure}[H]
\centering\includegraphics[width=\linewidth]{fig_convergence.pdf}
\caption{Reported identification traces in their original step coordinates. LIBERO
shows the median and interquartile range over 20 evaluation episodes. ALOHA and GR1
show identification-episode traces, followed by separate held-correction evaluations.
GR1 holds the median of reach-window estimates from three identification episodes.
Dashed lines mark injected offsets; these plots do not report execution-recovery time.}
\label{fig:convergence}
\end{figure}
Trace-level timestamps and exact fault-onset, gate, and correction-application events
are not available. Accordingly, the horizontal axes are not converted into seconds,
and inferred hold boundaries are not added. Nominal interface rates (50 Hz for ALOHA
and 20 Hz for GR1) alone do not establish the update or policy-call rate represented
by each logged step. Adapter latency is now measured over the eight healthy-control runs:
$.17$ ms median, $.46$ ms at the 99th percentile, $.77$ ms maximum per
control step, against a policy call near $900$ ms and a simulator step of
30--40 ms on the evaluation GPU. Synchronized sensor-to-command delay,
deadline misses, and a predefined sustained executed-action error threshold
for recovery remain unmeasured; the traces show parameter identification
behavior, not real-time guarantees.

\subsection{Measured execution certificates, estimator recursions, and the $r_y$ deficit}
\label{app:measured-certificates}

The theory above is conditional on constants the paper originally left open.
Each was preregistered with a refutation condition and then measured; the
outcomes split cleanly between the estimator, where the recursions are
predictive, and the plant, where the certificate the tube needs does not
exist on the primary robot.

\paragraph{Same-command contraction.}
Eq.~\eqref{eq:sampled-contraction} was measured directly: a simulator state
is saved mid-episode, perturbed, and the \emph{identical} recorded command
sequence is replayed from both states. Restoration is exact to $10^{-14}$
once the snapshot includes the solver warm start, actuator state, and the
gripper's accumulated target. Over 480 perturbations (three magnitudes, six
directions, four save points) on the LIBERO Panda under \code{OSC\_POSE},
the per-step gap ratio is $1.00$ in every preregistered metric -- joint,
joint-and-velocity, end-effector, and end-effector in the tracking metric of
the paragraph below -- with $47\%$ of free-space steps and $54\%$ of contact
steps at or above one. The controller sets each goal relative to the current
pose, so a joint-space perturbation neither decays nor grows. Per the
registered rule no further metric was tried. Contact raises joint-space
expansion and damps end-effector gaps ($32\%$ of steps above one, against
$50\%$ in free space). Measured input sensitivity is $.5$--$.6$\,cm per unit
command on $x,y,z$, a factor $2.3$--$2.9$ below the FIR's first tap.

\paragraph{Servo certificates.}
Per-joint first-order fits (leave-one-episode-out, held-out increment $R^2$)
certify the GR1's seven right-arm joints and waist yaw, poles $.72$--$.77$ at
20\,Hz with held-out $R^2$ $.74$--$.95$; the idle left arm and the hands fall
outside. ALOHA at 50\,Hz refutes the route on most joints -- only two arm
joints and the continuous gripper certify -- because its $20$\,ms increments
are noise-sized against the command--position gap. The Panda certifies
per-axis \emph{tracking} ($\lambda=.71,.67,.82,.83,.93,.61$ on
$x,y,z,r_x,r_y,r_z$; held-out $R^2$ $.84$--$.99$, stable to $\pm.02$ across
three initial states), which is a response model and not the same-command
contraction the tube requires; supplying it as the metric above still gives
$1.00$.

\paragraph{Tube and small-gain, propagated with measured constants.}
With $\lambda\approx1$, $L=.11$ and $\eta=.18$ from these measurements, the
tube of Theorem~\ref{thm:tube} covers the measured execution error on
$100\%$ of steps of 20 headline episodes with no violation to classify --
and is vacuous: it never forgets, growing by $\eta$ per step to a median
$98\times$ the largest error in its own episode. The small-gain condition of
Proposition~\ref{prop:smallgain} fails for the same reason and not for the
one it is about: fitted coupling constants are small ($\epsilon_0=.0008$,
$k_E=.012$, $k_X=.009$), but $a=1-\lambda=9\times10^{-5}$ gives
$ac=7\times10^{-6}$ against $bk_X=9.6\times10^{-4}$. On this plant the
composed certificate is unavailable, and the repair claim rests on the
paired counts.

\paragraph{Estimator recursions.}
On the same telemetry the estimator-side predictions hold. The envelope of
Proposition~\ref{thm:estimator}, propagated from the logged residual, gate
and attenuation, covers the measured selected-coordinate error on $100\%$ of
steps. Fixed points predicted from the \emph{logged residual distribution}
$\mathbb E[s_kz_k]$ match the stored ablation settles within $20\%$ on all
nine $(\rho,\text{channel})$ pairs. Under the ramp profile the estimate's lag
($.016$) sits inside $\nu/\underline\alpha=.021$. The registered onset
statistic fails because the plant's step noise is the size of the $.01$ step
it must resolve, but the episode-mean response follows the FIR partial sums
on $r_x$ and $r_z$ (plateau ratios $1.01$ and $.92$) and overshoots by
$3.0\times$ on $r_y$. The law's own analysis is therefore predictive where
the plant's is not, and the attenuation/innovation distinction appears
exactly as derived:

\begin{center}\small
\begin{tabular}{lrrr}
\toprule
Settled estimate, fraction of the $+.05$ truth & $r_x$ & $r_y$ & $r_z$\\
\midrule
Attenuation law (deployed) & $.87$ & $.41$ & $.88$\\
Innovation law, same scenarios & $.95$ & $.45$ & $.98$\\
\bottomrule
\end{tabular}
\end{center}
Removing the attenuation removes the $r_x,r_z$ shortfall, as
Proposition~\ref{prop:app-legacy} predicts, and leaves $r_y$ where it was.

\paragraph{The $r_y$ deficit is a calibration mismatch, not excitation or model order.}
Excitation and sensitivity are ruled out ($M$'s $r_y$ column is $.25$--$.28$
at three initial states, and $r_y$ is commanded as much as $r_z$). A
registered ARX(1) intervention, fitted on the same healthy log, repairs the
headline spatial cell $9/20\to20/20$ but settles $r_y$ at $32\%$, below the
FIR's $41\%$ and the registered $60\%$ floor: \emph{refuted}. The cause is
visible across plants. Every model fitted to the closed-loop healthy log
recovers nearly the same $r_y$ DC gain -- $.103$ (FIR $K{=}6$), $.118$
($K{=}20$), $.099$ (ARX) -- against the open-loop probe's $.276$, and the
onset residual plateaus at $.30$ per unit offset, confirming the probe. The
correction is executed at the probed gain but predicted at the fitted one, so
the applied estimate feeds back into the residual and the fixed point sits
low by that ratio:

\begin{center}\small
\begin{tabular}{lrrr}
\toprule
 & $r_x$ & $r_y$ & $r_z$\\
\midrule
Predicted ratio (fitted DC gain)/(probed $M$) & $.90$ & $.36$--$.43$ & $1.02$\\
Measured settled fraction, across laws & $.87$--$.95$ & $.32$--$.45$ & $.88$--$.98$\\
\bottomrule
\end{tabular}
\end{center}
A $2.7\times$ mismatch between closed-loop identification and the open-loop
probe, on the channel with the slowest tracking pole ($.93$), therefore
accounts for the shortfall on all three rotation channels. Constraining the
fitted DC gain to the probed $M$, or identifying from open-loop segments, is
the implied intervention; it is not registered and was not run.

\paragraph{Seed-pinned pairing.}
Pinning one sampler key per policy call makes two runs' frozen arms agree on
every outcome ($80/80$ and $20/20$ pairs) and execute identical commands to
the last step in $61/80$ and $17/20$ episodes, the remainder diverging at a
policy-call boundary from GPU nondeterminism; two adaptive arms under
different observers diverge at the first update, step $10$--$11$. Paired
comparisons within a pinned design therefore carry no sampling noise from
the policy's draws. Absolute levels do not transfer across designs: one key
per call is a single fixed noise realization rather than the policy's
average, and the pinned method scores $5/40$ on the long-horizon suite where
the unpinned cell scores $15/40$. Only within-design paired contrasts are
compared here.

\subsection{Hardware diagnostics and sensing limitations}

The X2 humanoid record contains 134 physical rollouts with stiffness, damping, and
feedforward scales from $0.1$ to $1.3$. The intervention is poorly visible in joint
position and better reflected in effort, with reported correlation $r=.84$.
The offset-detection range is $0.5$--$1.3^\circ$, and load-bearing joint behavior
varies across runs. These are sensing diagnostics; no paired physical-robot repair
result is included.

Motion residuals can confound actuator and sensor biases. The checkpoint-specific
zero-state-input ablation retains visual feedback and does not establish perceptual
robustness. These diagnostic limitations delimit the tested physical interface.
